\begin{resumo}

A educação é um pilar fundamental para o desenvolvimento de um país, mas o cenário educacional brasileiro enfrenta desafios como a baixa motivação e o desengajamento dos alunos, refletidos em indicadores como o Ideb e o PISA. Metodologias de aprendizagem colaborativa, como a tutoria, e a aplicação de gamificação têm demonstrado grande potencial para mitigar esses problemas, aumentando o engajamento e a eficiência do processo de ensino-aprendizagem. O objetivo deste trabalho é a implementação de uma aplicação web que intermedia o contato entre tutores de uma área de conhecimento e pessoas interessadas nesta tutoria. Trata-se de uma pesquisa aplicada, que utiliza procedimentos de pesquisa bibliográfica e metodologias ágeis de desenvolvimento de software. Ao final deste trabalho, foi implementada uma simulação da arquitetura para verificar a viabilidade da proposta, que se mostrou satisfatória ao atender os objetivos e validar a integração entre as tecnologias escolhidas para o êxito do projeto proposto.

 \vspace{\onelineskip}
    
 \noindent
 \textbf{Palavras-chaves}: Tutoria; Sistema de Recomendação; Gamificação; Aplicação Web.
\end{resumo}
